\chapter{Архитектурное решение после дефектно-транспортного аудита}

Часть II завершилась не выбором готовой модели, а точным условием
переоткрытия: пространственная семейная связность и ориентированный
спинорный подъём должны следовать из одного родительского принципа. После
ретроспективного сопоставления десяти возможных продолжений с Томами I--V
и строгой ветвью RPFT настоящее решение фиксирует один основной маршрут,
несколько вспомогательных языков и закрытые обходные пути.

Документ имеет статус архитектурного решения, а не новой физической
теории. Он определяет, какую гипотезу разрешено проверять следующей, но не
считает её истинной до прохождения перечисленных ниже гейтов.

\section{Принятое решение}

Основным кандидатом выбирается
\begin{equation}
 \boxed{
 \begin{gathered}
 \text{локальный }SU(2)_F\text{-подъём семейного переноса}
 \\
 +\;\text{градуированная суперсвязность Мориты}
 \end{gathered}}
 \label{eq:v5-decision-main-route}
\end{equation}
Индекс $F$ подчёркивает семейное происхождение конструкции. Он не означает
автоматического введения новой наблюдаемой низкоэнергетической
калибровочной группы.

Один локальный элемент
\begin{equation}
 U_{xy}\in SU(2)_F
\end{equation}
должен иметь два согласованных чтения:
\begin{enumerate}
 \item присоединённое действие
 \begin{equation}
  \operatorname{Ad}_{U_{xy}}\in SO(3)_F
 \end{equation}
 переносит семейную проекторную ось и в непрерывном пределе даёт
 пространственную связность конечной энергии;
 \item фундаментальное действие $U_{xy}$ переносит спинорный подъём,
 сохраняет потерянный проектором знак и допускает массовый оператор типа
 $\mathbf n\cdot\boldsymbol\sigma$ с единичным индексом.
\end{enumerate}

Связность, проектор, массовый оператор и физическое чтение должны входить
в одну градуированную суперсвязность
\begin{equation}
 \mathbb A=\nabla_F+\mathcal Q(P,\mathcal L)+\mathcal D_{\rm phys},
 \label{eq:v5-decision-superconnection}
\end{equation}
где $\nabla_F$ является связностью семейного перехода, $\mathcal L$ ---
данными двойного подъёма, а $\mathcal D_{\rm phys}$ --- уже существующим
стандарт-модельным чтением на $H_{15}$. Формула
\eqref{eq:v5-decision-superconnection} задаёт тип искомого объекта, но не
постулирует его конкретное действие.

\section{Почему выбран именно этот путь}

Выбор определяется не внешней привлекательностью монопольной модели, а
максимальным наследованием уже построенных результатов проекта.

\begin{enumerate}
 \item Том I требует, чтобы частица была устойчивым режимом единого закона
 допустимых переходов, а не отдельным добавленным объектом.
 \item Строгая RPFT уже использует двойное накрытие
 $S^3\to\mathbb{RP}^3$, спин-структуры и голономии как носители
 различий между векторным и спинорным чтениями.
 \item Том II запрещает независимые секторные коэффициенты и требует
 фиксировать нормировку до сравнения с наблюдаемыми.
 \item Том III требует одного действия и одной меры для нескольких
 чувствительных к нормировке секторов.
 \item Том IV уже построил семейный $SO(3)$, ранг-один проектор,
 тетраэдрическую остаточную структуру, точную голономию и условную
 семейную связность с одномерным ядром.
 \item Том V переосмыслил $M_{20\times15}(\mathbb C)$ как пространство
 переходов, построил локальный унитарный шаг, выделил цикл порядка четыре
 и доказал необходимость одновременно пространственной связности и
 спинорного подъёма.
\end{enumerate}

Следовательно, маршрут~\eqref{eq:v5-decision-main-route} добавляет не два
независимых поля, а проверяет, являются ли уже найденная семейная
$SO(3)$-связность и отсутствующий спинорный знак двумя представлениями
одного $SU(2)_F$-переноса.

\section{Распределение десяти рассмотренных идей}

\subsection{Компоненты основного маршрута}

\begin{enumerate}
 \item Семейный $SU(2)$-монополь принимается не как готовая модель, а как
 непрерывный предел локальной связности и проекторного дефекта.
 \item Дискретный $SU(2)$-автомат принимается как микроскопический язык
 перехода и возможный источник общей голономии.
 \item Суперсвязность принимается как контейнер, который обязан объединить
 связность, проектор и фермионную массу в одном кривизном объекте.
\end{enumerate}

\subsection{Вспомогательные структуры}

\begin{enumerate}
 \item $CP^1$-спинор допускается только как локальная координатизация
 фундаментального $SU(2)_F$-подъёма. Он не вводится самостоятельным полем.
 \item Структура типа $\operatorname{Spin}^h$ допускается как глобальный
 критерий существования совместного пространственно-семейного спинорного
 расслоения. Она не заменяет действие.
 \item Индуцированная фермионная связь допускается только после построения
 классического родительского оператора, как возможный источник конечных
 квантовых коэффициентов.
\end{enumerate}

\subsection{Отложенные и закрытые самостоятельные маршруты}

\begin{enumerate}
 \item Обычные внутренние флуктуации не переоткрываются: малая алгебра не
 видит семейное направление, а полная возвращает произвольную юкавскую
 свободу.
 \item Дефект Алисы не выбирается основным маршрутом, поскольку рискует
 связать точечный объект с физической струной или поверхностью разреза.
 \item Новое измерение или дополнительный круг запрещены до разрешения
 прежней развилки base--$K$/internal--$K$ и риска двойного KK-счёта.
 \item Член Скирма не используется как отдельный ремонт: он может
 стабилизировать размер текстуры, но не выводит спинорный знак и возвращает
 новый коэффициент.
\end{enumerate}

\section{Что означает подъём без скрытого расширения модели}

Недопустимо просто объявить новую калибровочную группу $SU(2)_F$ с новым
набором бозонов и свободной константой связи. Допустимы только два исхода.

\begin{enumerate}
 \item $SU(2)_F$ является структурной группой уже существующего семейного
 соответствия. Его связность выводится из проектора, метрики и правила
 параллельного переноса, а неприсоединённые к наблюдаемому сектору моды
 являются вспомогательными, топологическими, связанными или получают щель
 из того же механизма блокировки.
 \item Если такая интерпретация невозможна, $SU(2)_F$ является новым
 физическим сектором. Тогда маршрут отклоняется в текущей версии и может
 рассматриваться только как новая модель с повторным аудитом аномалий,
 спектра, меры и феноменологии.
\end{enumerate}

Тем самым выражение ``подъём $SO(3)_F$ к $SU(2)_F$'' означает прежде всего
проверку существования ассоциированного спинорного модуля, а не разрешение
добавить новые частицы.

\section{Первый обязательный гейт}

До построения монопольного профиля, потенциала или численного солитона
необходимо решить чистую представительную задачу.

Пусть $\mathcal E_{15}$ обозначает существующий наблюдаемый модуль, а
$\mathcal P_F\to X$ --- кандидат в главное $SU(2)_F$-расслоение. Требуется
проверить существование ассоциированного модуля $\mathcal S_F$ и
отображений
\begin{equation}
 \mathcal S_F\otimes\mathcal S_F^*\longrightarrow
 \operatorname{End}(\mathcal E_{15}),
 \qquad
 \operatorname{Ad}(\mathcal P_F)\longrightarrow
 \operatorname{End}(\mathcal E_{15}),
 \label{eq:v5-decision-representation-maps}
\end{equation}
таких, что одновременно выполнены следующие требования:

\begin{enumerate}
 \item присоединённое представление воспроизводит уже построенный семейный
 $SO(3)$-триплет и его проектор;
 \item центральный элемент $-1\in SU(2)_F$ действует как требуемый
 спинорный знак, но не создаёт нового наблюдаемого заряда;
 \item стандарт-модельные квантовые числа $H_{15}$ не изменяются;
 \item KO6, вещественная структура, градуировка и правило первого порядка
 сохраняются;
 \item не добавляются новые фермионные состояния;
 \item глобальные и смешанные аномалии отсутствуют;
 \item существующий след $M_{35}$ однозначно нормирует оба представления.
\end{enumerate}

\begin{nogocriterion}[Представительный стоп-критерий]
Если фундаментальный $SU(2)_F$-подъём нельзя реализовать на существующем
$H_{15}/M_{35}$ без нового фермионного модуля, независимой константы связи
или новой наблюдаемой калибровочной группы, основной маршрут закрывается до
записи любого монопольного действия.
\end{nogocriterion}

\section{Последовательность дальнейшей работы}

Только после прохождения представительного гейта разрешается следующая
цепочка.

\begin{enumerate}
 \item \textbf{Глобальный подъём.} Построить главное $SU(2)_F$-расслоение,
 его отношение к $SO(3)_F$, корневой линии порядка четыре и возможной
 структуре $\operatorname{Spin}^h$.
 \item \textbf{Геометрическая энергия.} Вывести из одной кривизны
 $F_A$, $D_AP$ и потенциальный член; проверить конечность энергии
 точечного дефекта без независимого члена Скирма.
 \item \textbf{Индексный оператор.} Получить
 $Q_{1/2}(P,\mathcal L)$ из того же подъёма и доказать единичный индекс,
 локализацию и спектральную щель.
 \item \textbf{Суперсвязностное действие.} Проверить, что единый след и
 одна мера фиксируют относительные веса калибровочной, проекторной и
 фермионной частей. Обычная сумма матриц Грама не считается успехом;
 требуемая ориентация должна следовать из относительной кривизны, а не из
 вставленного знака.
 \item \textbf{Физическое чтение.} Только затем применять уже построенную
 лестницу $300\to45\to15\to2\to1$ и сравнивать с наблюдаемыми.
\end{enumerate}

\section{Условия успеха и отказа}

Маршрут считается успешным лишь при одновременном выполнении:
\begin{equation}
 \boxed{
 \begin{gathered}
 \text{один }SU(2)_F\text{-перенос},
 \qquad
 E_{\rm defect}<\infty,
 \\
 \operatorname{ind}D_{\rm defect}=1,
 \qquad
 \text{нет новых фермионов},
 \\
 \text{нет независимых секторных весов},
 \qquad
 \text{один родительский след}.
 \end{gathered}}
\end{equation}

Маршрут закрывается, если любой из двух ключевых результатов --- конечная
энергия или единичный индекс --- требует отдельного поля, отдельной меры
или коэффициента, не фиксированного до сравнения с физикой.

\section{Финальный статус решения}

\begin{protocol}
\begin{enumerate}
 \item основной путь: \textbf{локальный $SU(2)_F$-подъём семейного переноса};
 \item родительский язык: \textbf{градуированная суперсвязность Мориты};
 \item $CP^1$ и $\operatorname{Spin}^h$: \textbf{вспомогательные чтения};
 \item обычные внутренние флуктуации: \textbf{не переоткрывать};
 \item Алиса, дополнительный круг и Скирм:
 \textbf{не использовать как самостоятельный ремонт};
 \item первый тест: \textbf{представительный подъём на $H_{15}/M_{35}$};
 \item монопольная динамика до первого теста: \textbf{запрещена};
 \item физическое замыкание: \textbf{не объявлено};
 \item новая версия программы: \textbf{не открыта до прохождения первого гейта}.
\end{enumerate}
\end{protocol}

Таким образом, принято не решение ``добавить $SU(2)$'', а более узкое
решение: проверить, был ли семейный $SO(3)$ всё время присоединённой тенью
одного уже допустимого спинорного переноса. Если ответ отрицателен, ветвь
закрывается дёшево и до введения новой феноменологии. Если положителен,
проект впервые получает один кандидат, способный связать конечную энергию,
двойной подъём и единичный индекс в одной архитектуре.

\section{Последующая проверка решения}

Обязательный представительный гейт выполнен в следующей главе и дал
отрицательный ответ. Текущий семейный модуль содержит только целоспиновые
представления, поэтому центр $SU(2)$ действует тривиально. KO6-удвоение,
слабый дублет и ориентационный множитель переноса не дают требуемый
семейный спинор без изменения градуировки, квантовых чисел или бюджета
модулей. Поэтому выбранный маршрут закрыт в версии V до монопольной
динамики; его возможное $\operatorname{Spin}^h$-продолжение относится уже
к отдельно объявляемой версии.

Последующий аудит проверил и эту возможность на минимальном бюджете.
Одна ориентационная $\mathbb C^2$ не может нести два независимых
коммутирующих действия --- пространственный спин и семейный полуцелый
изоспин. Настоящий $\operatorname{Spin}^h$-модуль требует как минимум
$\mathbb C^2_{\rm spin}\otimes\mathbb C^2_F$, то есть нового внутреннего
дублета и нового следа. Поэтому условие переоткрытия без расширения
моделей также не выполнено.

\section{Ретроспективная поправка после чтения раннего источника}

Позднейший аудит раннего текста ``Хабр 2'' обнаружил маршрут, не
попадающий под два предыдущих запрета. Полуцелое чтение можно искать не в
новом внутреннем $SU(2)_F$-дублете, а в хопфовой комплексной линии на
спинорном накрытии $S^3$. Её первый класс Черна равен единице, а
сопряжённые линии $L/L^*$ могут сопровождать направленные стрелки
$E/E^*$. Этот вариант не отменяет запреты $SU(2)_F$ и
$\operatorname{Spin}^h$, поскольку не вводит второй независимый
фундаментальный дублет. Он условно переоткрывает только вопрос о
функториальном происхождении ориентации стрелки.

Следующий гейт вывел это происхождение из двухугловой градуировки
$\Gamma_{\rm link}=p_{20}-p_{15}$. Она назначает степени $+1/-1$ стрелкам
$E/E^*$, а потому функториально связывает их с $L/L^*$. Неединственность
прежней трёхуровневой высоты при этом не отменяется: новый знак возникает
между неравноранговыми углами связывающей алгебры. Топологический мост
считается пройденным; открытым остаётся единый функционал энергии и
индекса.

Совместный энергетико-индексный гейт уточнил границу результата. Чистая
абелева хопфова связность не экранирует проекторный градиент и имеет
сингулярную энергию ядра. Однако тот же класс $c_1=1$ допускает гладкое
$SO(3)$-завершение типа Богомольного--Прасада--Соммерфилда: одна пара
$(A,\Phi)$ имеет конечную энергию и индекс Каллиаса один. Значит,
совместный механизм математически непротиворечив. Невыведенным остаётся
не его существование, а происхождение глобальной пространственной
суперсвязности, относительной жёсткости и масштабов из следа $M_{35}$.

Последующая проверка показала, что один след эту задачу не решает. Он не
порождает пространственную производную и не выбирает
$\mathfrak{so}(3)\subset\mathfrak u(35)$. Более того, линия с $c_1=1$ не
продолжается через ядро как одиночная линия: гладкий монополь требует
пары $L\oplus L^*$ и комплексно-линейного смешивания ветвей. Текущий
семейный модуль имеет только целые спины и даёт классы Черна $0$ или
чётные. Поэтому BPS-конструкция остаётся внешним математическим образцом.
Последняя внутренняя лазейка состоит в проверке, может ли требуемое
смешивание дать нечётная компонента уже существующей суперсвязности, а не
новый фиксированный семейный дублет.

Литературный аудит уточнил и эту постановку. В проверенной теории
суперсвязностей дефект естественно задаётся множеством потери обратимости
нечётной матричной массы. Свежие реализации 2026 года используют тот же
язык для произвольного комплексного матричного поля и локализации
топологического характера около его нулевого множества. Однако
естественный нечётный угол текущего родителя имеет размер $20\times15$ и
нигде не обратим: его коядро имеет размер не меньше пяти. Поэтому следующий
стоп-тест должен сначала вывести квадратный оператор $15\times15$ либо
$300\times300$ без ручного удаления или удвоения состояний. Только после
этого допустимы индекс и энергия.

Равногранный стоп-тест выполнен с дополнительной самопроверкой. Прямой
оператор $20\times15$ оставляет пять нулевых каналов на бесконечности и
имеет постоянный индекс $-5$; падение ранга на один добавляет два нуля,
но индекс не меняет. Канонического ранга-пятнадцать угла в $H_{20}$ нет,
пара $E/E^*$ связана антилинейной операцией и даёт скрытое удвоение до
600 состояний, а квадрат $150\times150$ потребовал бы универсального
частично-античастичного спаривания, несовместимого с заряженным $H_{15}$.
При этом разность рангов даёт точную суперразмерность
$\tau_{35}(p_{20}-p_{15})=1/7$. Она сохраняется как новый структурный
инвариант, но не является локализованным индексом. Внутренняя хопфова
ветвь версии V тем самым исчерпана.

Последующая типизационная сверка уточнила этот вывод. Голономный проектор
$P_0$ имеет ранг один в семейном триплете, поэтому его нельзя прямо
отождествлять с проектором ранга пять в $M_{35}$. Однако после помещения
обоих чтений в сравнительную алгебру
$M_{35}\otimes M_3\simeq M_{105}$ возникает точное стабильное тождество
\begin{equation}
 (20-15)\cdot3=15\cdot1,
 \qquad
 ([p_{20}]-[p_{15}])\otimes[I_3]=[p_{15}\otimes P_0]
 \quad\text{в }K_0(M_{105}).
\end{equation}
Тем самым повторение $1/7$ является не случайным числом, а общей
следовой двойственностью. Она ещё не задаёт частичную изометрию: для
положительного представителя требуется невыведенный проектор ранга пять
в $H_{20}$. Финальная заморозка поэтому откладывается на один строгий
тест граничной трансгрессии этого класса.

Литературная проверка локализовала этот тест до расширения Тёплица. Для
хопфовой граничной унитарной функции $u_H$ обёртывания один и
$q_0=p_{15}\otimes P_0$ естественная коэффициентная карта имеет вид
\begin{equation}
 \partial_{\mathcal T}([u_H]\boxtimes[q_0])=\pm[q_0].
\end{equation}
Её образ имеет ранг пятнадцать в $M_{105}$ и нормированный вес $1/7$.
Это канонический $K_0$-класс после фиксации расширения, но сам класс
расширения, поляризация Харди и вещественный $KKO$-подъём ещё не выведены
из $M_{35}$. Следующий тест должен построить эту точную
последовательность из хопфово-моритовских данных, а не вводить её как
внешний аналитический сектор.

Явная реализация этой карты использует односторонний сдвиг $S$ на
$H^2(S^1)$ и коэффициентный проектор $q_0$. Оператор
$S\otimes q_0$ имеет компактный дефект
\begin{equation}
 (I-SS^*)\otimes q_0
 =|e_0\rangle\langle e_0|\otimes q_0
\end{equation}
ранга пятнадцать и веса $1/7$. Сопряжённая ориентация даёт индекс
противоположного знака, поэтому KO6-пара имеет индексы $-15/+15$ и
нулевой обычный суммарный индекс. Комплексный классовый мост тем самым
замкнут; пространство Харди остаётся аналитическим представителем, а
полный вещественный $KKO$-цикл и физический родительский оператор ещё не
построены.

Вещественный подъём этой конструкции даёт более сильное ограничение.
При знаках текущего KO6
$J^2=1$, $JF=FJ$, $J\gamma=-\gamma J$ оператор $J$ задаёт биекцию между
$\ker T$ и $\ker T^*$. Поэтому обычный индекс полного вещественного
пакета обязан быть равен нулю. Явная сбалансированная пара имеет индексы
$-15/+15$, суммарный компактный ранг тридцать на удвоенном носителе 210
и по-прежнему вес $1/7$. Ненулевой индекс сохраняется только у выбранной
комплексной ориентации; для полного пакета требуется отдельная
классификация Real/KR-класса после фиксации вещественной алгебры,
инволюции и Clifford-степени.

Эта классификация теперь выполнена. Поточечное сопряжение одной ветви
дало бы неподвижную алгебру $M_{105}(\mathbb R)$ и нулевую группу
$KO_6$, но оно не реализует проектный обмен $L\leftrightarrow L^*$.
Текущий $J$ выбирает обменную инволюцию
\begin{equation}
 \rho(a,b)=(\overline b,\overline a)
 \quad\text{на}\quad
 M_{105}(\mathbb C)\oplus M_{105}(\mathbb C).
\end{equation}
Её неподвижная вещественная алгебра изоморфна
$M_{105}(\mathbb C)_{\mathbb R}$, а потому
\begin{equation}
 KO_6\bigl(M_{105}(\mathbb C)_{\mathbb R}\bigr)\simeq\mathbb Z.
\end{equation}
Следовательно, сокращение комплексных индексов $-15/+15$ не доказывает
нулевой вещественный класс. Кандидат абсолютной величины пятнадцать и
веса $1/7$ сохраняется, но его ещё нельзя считать доказанным индексом:
следующий гейт обязан построить явный $Cl_{0,6}$-линейный фредгольмов
цикл и проверить отображение его класса в две комплексные ориентации.

Обязательная ретроспектива уточнила этот план. Тома III--V уже содержат
универсальное KO6-удвоение с противоположными градуировками и точным
правилом $JdJ^{-1}=d^\dagger$. Следовательно, повторное построение
конечномерного particle--antiparticle модуля не требуется и могло бы
ошибочно расширить физический спектр. Архивная строгая RPFT-ветвь также
не даёт недостающего индекса: для семейства голономий $0\to\pi$ в ней
самой получен нулевой спектральный поток.

Поэтому ближайший гейт сужается до отображения сравнения
\begin{equation}
 KO_6(\mathbb C_{\mathbb R})\longrightarrow
 K_0(\mathbb C\oplus\mathbb C).
\end{equation}
Необходимо проверить, переходит ли генератор в антидиагональную пару
$(-1,+1)$ с точностью до общего знака. Только после этого умножение на
$[q_0]$ ранга пятнадцать докажет пару $(-15,+15)$.

Этот тест теперь выполнен по явной таблице объединённой $K$-теории.
Для комплексных чисел как вещественной алгебры
\begin{equation}
 c_6(n)=(-n,n),
 \qquad
 r_6(a,b)=-a+b.
\end{equation}
Следовательно, $c_6$ инъективно,
$r_6c_6(n)=2n$ и
\begin{equation}
 c_6(15)=(-15,+15).
\end{equation}
Построенный обменный цикл поэтому представляет единственный класс 15 в
группе
\[
 KO_6(M_{105}(\mathbb C)_{\mathbb R}).
\]
Классификационный разрыв закрыт; открытым остаётся вывод неограниченного
родительского оператора и физического действия.

Аналитический неограниченный представитель теперь получен без внешнего
пространства Харди. На $\ell^2(\mathbb Z)$ оператор числа
$Ne_n=ne_n$ имеет спектральный проектор
$P=\chi_{[0,\infty)}(N)$, а
\begin{equation}
 PUP=S,
 \qquad
 PU^*P=S^*.
\end{equation}
Следовательно, поляризация, компактный дефект и индексы выводятся из
одного extension-цикла. Но этот неограниченный модуль имеет степень один.
При переводе из обозначения $KKO_1$ для extension-класса в ковариантную
точную последовательность $KO_n$ знак степени необходимо обратить:
boundary имеет вид $KO_n\to KO_{n-1}$. Поэтому для факторизации
доказанного класса $KO_6$ требуется входной обменный Real-символ степени
семь, поскольку $7-1=6\pmod8$. Сам класс пятнадцать при этом остаётся
доказанным.

Исправленная степень допускает каноническую факторизацию. Пусть
$\kappa_{15}=15\in KO_6(M_{105}(\mathbb C)_{\mathbb R})$ --- доказанный
обменный класс, а $[u_{\mathbb R}]\in KO_1(C^*_{\mathbb R}(\mathbb Z))$
--- winding-класс двустороннего сдвига. Тогда
\begin{equation}
 \xi_{15}=[u_{\mathbb R}]\boxtimes\kappa_{15}\in KO_7
\end{equation}
и модульность граничной карты даёт
\begin{equation}
 \partial_7(\xi_{15})=\pm\kappa_{15}.
\end{equation}
При согласованной ориентации комплексификация равна $(-15,+15)$.
Классовая неограниченная факторизация тем самым закрыта; явная единичная
матричная модель $KO_7$ и физическое действие остаются открытыми.

Единичная модель теперь также построена:
\begin{equation}
 V_{15}(z)=\bigl(zq_0+1-q_0,
 z^{-1}\overline q_0+1-\overline q_0\bigr).
\end{equation}
Она удовлетворяет $V_{15}^{\tau}=V_{15}$, имеет winding-пару
$(15,-15)$ и Toeplitz boundary $(-15,+15)$. Наивная формула
$pu+(1-p)u^*$ исключена, поскольку удваивает класс до тридцати.
Топологическая операторная факторизация закрыта; открытым остаётся
вариационное физическое действие.

Минимальный функционал, следующий из уже построенных данных, имеет вид
\begin{equation}
 S_{\rm loop}(V)=\tau_{210}([N,V]^*[N,V]).
\end{equation}
При фиксированном winding он выбирает $V_{15}$ и даёт значение $1/7$.
Но при изменении радиуса ведёт себя как $R^{-2}$ или, с мерой окружности,
как $R^{-1}$. Поэтому внутренняя петля вариационно устойчива в своём
топологическом секторе, тогда как конечный пространственный размер и
масса не стабилизированы.